\documentclass[12pt,a4paper]{article}
\usepackage[spanish,es-nodecimaldot]{babel}
\usepackage{fontspec}
% Instalar Times New Roman y compilar con XeLaTeX para usar la fuente solicitada.
\setmainfont{Times New Roman}
\usepackage[margin=2.3cm]{geometry}
\usepackage{longtable,array,booktabs,fvextra,enumitem,fancyhdr,needspace}
\usepackage[hidelinks,unicode]{hyperref}
\usepackage{xurl}
\setlength{\parindent}{0pt}
\setlength{\parskip}{6pt}
\setlength{\emergencystretch}{3em}
\setlist{nosep,leftmargin=1.5em,topsep=5pt}
\pagestyle{fancy}\fancyhf{}
\fancyhead[L]{\small IDENTITY ACCESS LAB}\fancyhead[R]{\small Manual de laboratorio}
\fancyfoot[C]{\thepage}
\setlength{\headheight}{15pt}
\DefineVerbatimEnvironment{Comandos}{Verbatim}{breaklines,breakanywhere,fontsize=\small,fontfamily=\rmdefault,frame=single,framesep=5pt}
\begin{document}
{\LARGE\bfseries Manual de implementación de IDENTITY ACCESS LAB\par}\vspace{6pt}
2 de octubre de 2026\par\hrule\vspace{10pt}
Fecha: 2 de octubre de 2026. Este paquete público explica cómo construir un laboratorio genérico. Todos los identificadores y direcciones de ejemplo son ficticios. Los pasos no significan que las plataformas ya estén configuradas.\par
\Needspace{5\baselineskip}\section*{Decisiones del montaje}
\begin{itemize}
\item NVR virtual autorizado: \textbf{Frigate}, software de código abierto. La configuración incluida usa la versión 0.17.1 para que todos trabajen con la misma base.
\item Cámaras: teléfonos Android y iPhone. La ruta Android usa IP Webcam y video MJPEG; la ruta iPhone usa IP Camera Lite y video RTSP/H.264. Las aplicaciones de teléfono no se presentan como software de código abierto; ese requisito se aplica al NVR.
\item Identidades digitales: \textbf{OneLogin obligatorio}, con inicio de sesión real y una conexión para crear y retirar cuentas de la aplicación.
\item Microsoft 365: administración de cuentas, licencias, roles, contraseña, MFA y registros en un entorno de práctica autorizado.
\item Servidor y NVR: comandos para Ubuntu 24.04 LTS de 64 bits. Si la laptop usa Windows, preparar una máquina virtual Ubuntu o conseguir una laptop Linux; no ejecutar comandos de Ubuntu en PowerShell.
\item Dos laptops disponibles por departamento: 14 equipos con funciones asignadas. Las demás apoyan desarrollo, documentación y pruebas.
\item RFID: modelo pendiente. El ejemplo incluido es únicamente para Arduino Uno + MFRC522. Identificar el equipo antes de cablear.
\end{itemize}
\Needspace{5\baselineskip}\section*{Archivos}
Las ocho guías están en \textnormal{Guias/\allowbreak{}}. Los diagramas están separados en \textnormal{Diagramas/\allowbreak{}}, en SVG y PDF, con fuentes LaTeX. El código y las configuraciones están en \textnormal{Implementacion/\allowbreak{}}. Los formularios y listas están en \textnormal{Plantillas/\allowbreak{}}. \textnormal{REFERENCIAS\_\allowbreak{}Y\_\allowbreak{}VIDEOS.pdf} reúne documentación y videos vinculados desde cada guía.\par
\Needspace{5\baselineskip}\section*{Orden de trabajo}
\begin{enumerate}
\item Dirección y RRHH corrigen la lista de integrantes y acuerdan responsables.
\item TI prepara red, equipos Ubuntu y conexión a internet.
\item Seguridad Física identifica el lector y monta cámaras y NVR con apoyo de TI.
\item TI prepara la aplicación. IAM configura OneLogin y conecta inicio de sesión y altas/bajas.
\item TI configura Microsoft 365; IAM coordina la segunda verificación.
\item SOC prepara monitoreo y consulta de registros.
\item Auditoría ejecuta pruebas con todos los equipos.
\item Se realiza el recorrido completo y se prepara la entrega.
\end{enumerate}
No esperen a terminar todo el RFID para iniciar red, cámaras y cuentas. La integración de OneLogin depende de que la aplicación y su dirección HTTPS funcionen primero.\par
\Needspace{5\baselineskip}\section*{Distribución propuesta}
31 puestos genéricos: Dirección 3, RRHH 3, TI 7, IAM 5, SOC 5, Seguridad Física 5 y Auditoría 3. Los puestos no representan personas identificadas. Las asignaciones y cualquier cambio se mantienen en un registro privado.\par
\Needspace{5\baselineskip}\section*{Equipos por área}
{\small\renewcommand{\arraystretch}{1.2}
\begin{longtable}{p{5.027cm}p{5.027cm}p{5.027cm}}
\toprule
\textbf{Área} & \textbf{Laptop 1} & \textbf{Laptop 2} \\ \midrule\endfirsthead
\textbf{Área} & \textbf{Laptop 1} & \textbf{Laptop 2} \\ \midrule\endhead
Dirección & Calendario, pendientes y control del grupo & Guion, presentación y copia de entregables \\
RRHH & Registro de empleados & Solicitudes de alta, cambio y baja \\
TI & Servidor de aplicación y conexión HTTPS & Red, Microsoft 365 y soporte \\
IAM & Consola OneLogin & Pruebas de usuarios y permisos \\
SOC & Cámaras desde el navegador & Registros, incidentes y evidencias \\
Seguridad Física & Microcontrolador y registro RFID & Frigate y almacenamiento de video \\
Auditoría & Lista de pruebas & Informe, referencias y evidencias \\
\bottomrule\end{longtable}}
Las funciones son asignaciones propuestas, no una exigencia de mantener todas las laptops encendidas a la vez. El NVR y servidor deben permanecer encendidos durante las pruebas continuas.\par
\Needspace{5\baselineskip}\section*{Direcciones de ejemplo}
{\small\renewcommand{\arraystretch}{1.2}
\begin{longtable}{p{7.760cm}p{7.760cm}}
\toprule
\textbf{Equipo} & \textbf{Dirección prevista} \\ \midrule\endfirsthead
\textbf{Equipo} & \textbf{Dirección prevista} \\ \midrule\endhead
Router/punto de acceso & 10.31.0.1 \\
Servidor TI & 10.31.0.10 \\
NVR & 10.31.0.20 \\
Laptop RFID & 10.31.0.21 \\
Android de entrada & 10.31.0.31 \\
iPhone interior & 10.31.0.32 \\
SOC & 10.31.0.40 y 10.31.0.41 \\
\bottomrule\end{longtable}}
Esta red es una propuesta. Si el router usa otra red, TI actualiza direcciones, diagramas y configuraciones. No asignar direcciones de ejemplo que no pertenezcan a la red real. Preferir reservas de dirección en el router.\par
\Needspace{5\baselineskip}\section*{Qué representa cada sistema}
\begin{itemize}
\item La tarjeta y el LED deciden el paso físico supervisado.
\item Frigate guarda video y permite consultarlo desde laptop y móvil.
\item OneLogin comprueba la identidad del usuario y aplica su política de MFA.
\item La aplicación comprueba qué acciones puede realizar el usuario.
\item La conexión SCIM transmite a la aplicación altas, cambios y bajas desde OneLogin. SCIM es el mecanismo de sincronización, no otro proveedor de identidades.
\item Microsoft 365 administra las cuentas corporativas del ejercicio. En este diseño no se federan ni sincronizan automáticamente Microsoft 365 y OneLogin. Se relacionan mediante el registro de RRHH y el mismo correo de laboratorio. Cambiar la contraseña en Microsoft 365 no cambia automáticamente la de OneLogin.
\end{itemize}
Se usan dos conectores en OneLogin: uno OIDC para iniciar sesión y uno SCIM Core para sincronizar usuarios. El conector SCIM incluye “SAML” en su nombre comercial, pero aquí solo se usa su función de aprovisionamiento; el inicio de sesión de la aplicación es OIDC.\par
\Needspace{5\baselineskip}\section*{Datos que hay que completar}
Dirección registra sistemas operativos, modelos del lector y placa, acceso a tenants, licencias, puerto/dirección de cada cámara y duración acordada de la prueba 24/7. La autorización del NVR virtual ya está confirmada. La adaptación específica de transceptores de video se debe dejar registrada por escrito para que el informe explique cómo se cubrió o sustituyó ese punto.\par
\Needspace{5\baselineskip}\section*{Uso de comandos y verificación}
\begin{itemize}
\item Copiar el paquete a las laptops designadas. Las rutas \textnormal{\textasciitilde{}/\allowbreak{}lab-runtime/\allowbreak{}} se refieren a la carpeta del proyecto en esas laptops.
\item Reemplazar todo texto \textnormal{REEMPLAZAR} antes de iniciar servicios.
\item No compartir \textnormal{.env}, claves privadas, contraseñas ni códigos de MFA en evidencias.
\item Ejecutar cada bloque en el equipo indicado y revisar el resultado antes del siguiente paso.
\item El código base tiene pruebas locales de permisos, cambio de rol, baja y vínculo de identidad. La conexión real con OneLogin, las cámaras y el RFID requiere sus cuentas y equipos; no se considera probada en este paquete.
\end{itemize}
Consultar referencias y videos (REFERENCIAS\_Y\_VIDEOS.pdf). Los videos son apoyo visual; las configuraciones versionadas y las fuentes oficiales resuelven diferencias de menús o versiones.\par
\end{document}
